% ============================================================
% 国赛格式 LaTeX 骨架（配套 math-modeling-paper 技能）
% 编译：tectonic main.tex（Tectonic 会自动下载 fandol 字体与宏包）
% ============================================================
\documentclass[11pt,a4paper,fontset=fandol]{ctexart}
\usepackage{amsmath,amssymb}
\usepackage{graphicx}
\usepackage{xcolor}
\usepackage{geometry}
\geometry{a4paper, top=2.5cm, bottom=2.5cm, left=2.5cm, right=2.5cm} % 左右必须对称
\usepackage{booktabs}
\usepackage{tabularx}
\newcolumntype{Y}{>{\raggedright\arraybackslash}X}   % 左对齐自适应列，防 underfull
\usepackage{adjustbox}
\usepackage{xurl}                                    % URL 可在斜杠处断行
\usepackage{needspace}                               % 防小节标题孤悬页底
\usepackage{enumitem}
\usepackage{tikz}
\usetikzlibrary{arrows.meta,positioning}
\usepackage{caption}
\captionsetup{font=small,labelfont=bf}
\usepackage{listings}                                % 代码附录
\lstset{basicstyle=\tiny\ttfamily, numbers=left, numberstyle=\tiny\color{gray},
  breaklines=true, frame=single, framerule=0.3pt, columns=flexible, language=Python}
\usepackage[numbers,super,sort&compress]{natbib}     % 上标引用
\widowpenalty=10000
\clubpenalty=10000

\usepackage[colorlinks=true, linkcolor=blue!60!black, citecolor=blue!60!black,
            urlcolor=blue!60!black, bookmarks=true, bookmarksnumbered=true,
            unicode=true]{hyperref}                  % 永远最后加载
\hypersetup{
    pdftitle={论文标题}, pdfauthor={作者}, pdfsubject={竞赛题号},
    pdfkeywords={关键词1, 关键词2}}

\newcommand{\um}{\,\mu\mathrm{m}}

\begin{document}

% ---------- 第 1 页：标题 + 摘要 + 关键词（国赛格式：无封面无目录） ----------
\begin{center}
    {\zihao{3}\heiti 论文标题（可两行）}
\end{center}

\begin{center}
    {\zihao{4}\heiti 摘\quad 要}
\end{center}

本文针对……问题，建立了……模型，设计了……算法。（总起段：对象-方法-结论一句话）

\textbf{针对问题一}，我们……（方法一句话）。……（推导要点）。得到……（数值结论）。

\textbf{针对问题二}，我们设计了……算法：……（一句话机理）。对附件 1、2 的结果为
$d=7.20\pm0.09\um$（95\% 置信区间 $[7.01,7.38]\um$）。
% ↑ 摘要里每个结论必须带数值，且与正文/表格完全一致

\textbf{针对问题三}，……判定 + 建模 + 数值 + 复核……

\vspace{0.6em}
\noindent{\heiti 关键词：}关键词一；关键词二；\mbox{最后一个关键词}  % \mbox 防拆行

\newpage

\section{问题重述}
% 用（1）（2）（3）列出每问要求，忠于原文措辞

\section{问题背景与研究现状}
% 行业背景一段 + 方法学四条主线各一段（极值法/频域法/包络法/全谱拟合法）+ 开源实践现状

\section{问题分析}
% 每问一小节：候选方案对比表（优点/缺点/本文取舍）——评审最看重"为什么这么做"

\section{模型假设}
\begin{enumerate}[label=\textbf{假设\arabic*：}, leftmargin=3.6em]
    \item ……。\emph{依据}：…… % 每条假设必须带物理依据
\end{enumerate}

\section{符号说明}
% booktabs 三线表：符号 | 含义 | 单位/取值

\section{问题一：模型的建立}
% 从物理基本关系推导；关键公式 \label 供后文 \eqref 引用

\section{问题二：算法设计与求解}
% 预处理 → 特征提取 → 主算法（TikZ 流程图，≤6 节点）→ 交叉验证方法×2
% → 结果汇总表 → 可靠性（双数据互证/Bootstrap/灵敏度/蒙特卡洛）

\section{问题三：……}
% 判定（定量判据）→ 建模 → 算法（全局优化+局部精修，收敛曲线图）→ 结果 → 对照 → 复核

\section{与开源实现的系统对照}
% 表：各实现 | 物理模型 | 算法 | 报告数值；讨论方法间系统离散的归因

\section{灵敏度与稳健性检验}\label{sec:sens}
% 对数微分传递公式 + 数值扰动表 + 蒙特卡洛直方图

\section{模型评价与推广}
% 优点（对应创新点）/ 不足（诚实量化方法间离散）/ 推广

\begin{thebibliography}{99}
    \raggedright
    \setlength{\itemsep}{0.15em}
    \small
    \bibitem{ref1} 作者. 题名[J]. 刊名, 年, 卷(期): 页码. DOI: xxx.
    % 每条文献写进论文前必须在 Crossref 核验过
\end{thebibliography}

\newpage
\appendix
\section{附录：Python 关键代码}
\begin{lstlisting}
# 核心算法代码（能跑通的最小集）
\end{lstlisting}

\end{document}
